\documentclass[11pt]{article}
\usepackage[margin=1in]{geometry}
\usepackage{amsmath,amssymb,amsthm,mathtools}
\usepackage{mathrsfs}
\usepackage[hidelinks]{hyperref}
\newtheorem{theorem}{Theorem}
\newtheorem{lemma}[theorem]{Lemma}
\newtheorem{proposition}[theorem]{Proposition}
\newcommand{\SU}{\operatorname{SU}}
\newcommand{\tr}{\operatorname{tr}}
\newcommand{\Rea}{\operatorname{Re}}
\newcommand{\dd}{\,\mathrm d}
\title{Exact SU(2) Wilson transfer operator and the finite-spatial-lattice time limit}
\author{}
\date{8 September 2026}
\begin{document}
\maketitle

\section{Retained action, coefficients, and conventions}

Fix $a>0$ and $g_{\rm YM}>0$. The spatial regulator is a fixed finite
oriented cubical graph with vertices $\mathsf V$, physical spatial edges
$\mathsf E$, and the full set $\mathsf P_s$ of its geometric spatial
plaquettes, each included once, with $|\mathsf E|\geq1$.
Open spatial boundary conditions or a
fixed periodic spatial lattice are both allowed. Reverse traversal uses
$U_e^{-1}$, not a second independent link. Write $N=|\mathsf E|$ and
\[
 \mathcal Q=\SU(2)^{\mathsf E},\qquad
 \mathcal H=L^2(\mathcal Q,\lambda),\qquad
 d\lambda=\prod_{e\in\mathsf E}dU_e,
\]
with probability Haar on each link. The spatial graph, $a$, and
$g_{\rm YM}$ stay fixed in every limit in this paper.

The original invariant Lie metric and continuum action are
\begin{equation}\label{eq:original-action}
 c(X,Y)=-\frac12\tr(XY),\qquad
 S(\mathscr A)=\frac1{2g_{\rm YM}^2}
 \int\sum_{\mu<\nu}c(F_{\mu\nu},F_{\mu\nu})\,d^4x.
\end{equation}
For time spacing $\epsilon>0$, retain exactly
\begin{equation}\label{eq:anisotropic}
 b_t(\epsilon)=\frac{a}{2g_{\rm YM}^2\epsilon},\qquad
 b_s(\epsilon)=\frac{\epsilon}{2g_{\rm YM}^2a}.
\end{equation}
For the frozen-curvature holonomy $e^{\delta F}$,
$2-\Rea\tr e^{\delta F}=\delta^2c(F,F)+O_F(\delta^4)$ on
$\mathfrak{su}(2)$, by the exponential series, tracelessness and the
paired eigenvalues of $F$. The exact coefficient checks are
\[
 \frac{b_t(a\epsilon)^2}{a^3\epsilon}
 =\frac1{2g_{\rm YM}^2},\qquad
 \frac{b_sa^4}{a^3\epsilon}=\frac1{2g_{\rm YM}^2}.
\]
These equalities record the prescribed classical action convention.
The transfer calculation below uses the full compact plaquette action,
not a curvature truncation, and proves no spatial continuum limit.

Choose $T_1,T_2,T_3\in\mathfrak{su}(2)$ with
$-2\tr(T_aT_b)=\delta_{ab}$, for example $T_a=i\sigma_a/2$ with the
Pauli matrices. Define
\[
 X_{e,a}f(U)=\left.\frac d{dt}\right|_{0}
 f(\ldots,e^{tT_a}U_e,\ldots),\qquad
 \Delta_e^T=\sum_{a=1}^3X_{e,a}^2.
\]
In the original metric, $c(T_a,T_b)=\delta_{ab}/4$, so the
$c$-orthonormal basis is $S_a=2T_a$, and
\begin{equation}\label{eq:metricdictionary}
 X^c_{e,a}=2X_{e,a},\qquad \Delta_e^c=4\Delta_e^T.
\end{equation}
The probability Haar convention changes no operator coefficient. If
$d\operatorname{vol}_c=\mathcal V_c d\lambda$ on this compact product,
$f\mapsto\sqrt{\mathcal V_c}f$ is the exact unitary change from
$L^2(d\operatorname{vol}_c)$ to $L^2(d\lambda)$ and intertwines the same
differential operators.

The original inverse-transport gauge law and the convention used in the
companion Hamiltonian derivation are related at the configuration level by
\begin{equation}\label{eq:gaugedictionary}
 (\beta_hU)_e=h_{s(e)}^{-1}U_eh_{t(e)},\qquad
 (\alpha_gU)_e=g_{s(e)}U_eg_{t(e)}^{-1},\qquad
 \alpha_gU=\beta_hU\ \text{when }g_v=h_v^{-1}\text{ for every }v.
\end{equation}
Inversion preserves Haar, so their vertex gauge averages are the same
operator. No link variable is inverted by this dictionary; only the
parameter describing the same vertex change is replaced by its inverse.

\section{The exact single-link character multiplier}

For $b>0$, let
\begin{equation}\label{eq:kernel}
 Z(b)=\int_{\SU(2)}e^{b\Rea\tr U}dU,\qquad
 k_b(U)=\frac{e^{b\Rea\tr U}}{Z(b)},\qquad
 (\mathcal K_b f)(U')=\int k_b(U'U^{-1})f(U)dU.
\end{equation}
Every $U\in\SU(2)$ has eigenvalues $e^{i\theta},e^{-i\theta}$ with
$0\leq\theta\leq\pi$. Probability Haar integration of a class function is
\[
 \int f(U)dU=\frac2\pi\int_0^\pi f(\theta)\sin^2\theta\,d\theta.
\]
For example, write $U=w_0I+i\sum_{a=1}^3w_a\sigma_a$ on the unit
three-sphere. Slicing at $w_0=\cos\theta$ gives density
$(2/\pi)\sqrt{1-w_0^2}dw_0$, which is the displayed formula.

Define the integer-order modified Bessel functions by their entire series
\begin{equation}\label{eq:Besselseries}
 I_d(2b)=\sum_{r=0}^{\infty}\frac{b^{2r+d}}{r!(r+d)!},
 \qquad d\in\mathbb Z_{\geq0}.
\end{equation}
They also satisfy
\begin{equation}\label{eq:Besselintegral}
 I_d(2b)=\frac1\pi\int_0^\pi e^{2b\cos\theta}\cos(d\theta)\,d\theta.
\end{equation}
To prove this, multiply the uniformly absolutely convergent series for
$e^{bz}$ and $e^{b/z}$ on $|z|=1$. The coefficient of $z^d$ is
\eqref{eq:Besselseries}, and Fourier orthogonality gives
\eqref{eq:Besselintegral}. For $d\geq1$, integration of the derivative
of $e^{2b\cos\theta}\sin(d\theta)$ over $[0,\pi]$ gives
\begin{equation}\label{eq:recurrence}
 I_{d-1}(2b)-I_{d+1}(2b)=\frac db I_d(2b).
\end{equation}
The endpoint term is zero, and
$2\sin\theta\sin(d\theta)=\cos((d-1)\theta)-\cos((d+1)\theta)$
fixes every coefficient in this recurrence.

The spin-$j$ irreducible representation has $j\in\frac12\mathbb Z_{\geq0}$,
dimension $d=2j+1$, and character
\[
 \chi_j(\theta)=\sum_{r=0}^{d-1}e^{i(d-1-2r)\theta}
 =\frac{\sin(d\theta)}{\sin\theta}.
\]
The finite sum is the character of the $(d-1)$st symmetric power of the
defining representation; its equality with the sine quotient is the
finite geometric-series identity, including endpoint values by
continuity. These are all irreducibles of $\SU(2)$, and their matrix
coefficients are complete by Peter--Weyl \cite{PeterWeyl}.

\begin{theorem}\label{thm:multiplier}
The exact normalization and the exact spin-$j$ multiplier of
\eqref{eq:kernel} are
\begin{equation}\label{eq:multiplier}
 Z(b)=\frac{I_1(2b)}b,\qquad
 \mathcal K_bD^j_{rs}=m_j(b)D^j_{rs},\qquad
 m_j(b)=\frac{I_{2j+1}(2b)}{I_1(2b)}.
\end{equation}
For every $j$ and $b>0$, $0<m_j(b)\leq1$, and $m_0(b)=1$.
The operator $\mathcal K_b$ is a positive self-adjoint contraction,
is injective, and is trace class.
\end{theorem}
\begin{proof}
The exact unnormalized character integral is
\begin{align*}
 \int e^{b\Rea\tr U}\chi_j(U)dU
 &=\frac2\pi\int_0^\pi e^{2b\cos\theta}
                      \sin\theta\sin(d\theta)d\theta\\
 &=I_{d-1}(2b)-I_{d+1}(2b)=\frac dbI_d(2b).
\end{align*}
Set $d=1$ to obtain $Z(b)$. Since $k_b$ is central, the matrix
$\int k_b(W)D^j(W^{-1})dW$ commutes with the representation and hence
is scalar. Its trace divided by $d$ is
$(dZ(b))^{-1}\int e^{b\Rea\tr W}\chi_j(W)dW=I_d/I_1$.
The substitution $U=W^{-1}U'$ in \eqref{eq:kernel} proves the asserted
matrix-coefficient eigenvalue. In particular the dimension factor in
the character integral is exactly canceled by division by $d$; no
dimension factor remains in the multiplier.

All terms in \eqref{eq:Besselseries} are positive for $b>0$, which
proves $m_j>0$. The integral expression for the normalized multiplier
and $|\chi_j(U)|\leq d$ prove $|m_j|\leq1$. The matrix coefficients
form an orthogonal complete basis, so these real positive eigenvalues
prove positivity, self-adjointness, contraction, and injectivity.
For fixed $b$, \eqref{eq:Besselseries} gives
$I_d(2b)\leq b^de^{b^2}/d!$, since $(r+d)!\geq d!$.
Thus
\[
 \tr\mathcal K_b=\sum_{d=1}^{\infty}d^2\frac{I_d(2b)}{I_1(2b)}<\infty,
\]
where the multiplicity is exactly $d^2$. The factorial majorant proves
the convergence and hence trace-classness.
\end{proof}

\section{Controlled small-time derivative of the full compact kernel}

\begin{lemma}\label{lem:asymptotic}
For each fixed spin $j$,
\begin{equation}\label{eq:large-b}
 m_j(b)=1-\frac{j(j+1)}b+O_j(b^{-2})\qquad(b\longrightarrow\infty).
\end{equation}
The estimate is proved from the exact compact integral and does not
replace that integral by a different finite-$b$ kernel.
\end{lemma}
\begin{proof}
Set $y=1-\cos\theta\in[0,2]$. Under the probability density $k_b dU$,
the density of $y$ is proportional to
$e^{-2by}y^{1/2}(2-y)^{1/2}dy$. Write
\[
 J_q(b)=\int_0^2e^{-2by}y^{q+1/2}(2-y)^{1/2}dy,
 \qquad q=0,1,2.
\]
On $[0,1]$, Taylor's theorem with a bounded second derivative gives
$\sqrt{2-y}=\sqrt2(1-y/4)+R(y)$, with $|R(y)|\leq C y^2$.
The part on $[1,2]$ is bounded by a constant times $e^{-2b}$.
The tails when extending the two leading integrals to $[0,\infty)$
have the same exponentially small order times a polynomial in $b$.
Finally, the error integral over $[0,1]$ is bounded by
$C\int_0^\infty e^{-2by}y^{q+5/2}dy$.
Using the exact substitution $z=2by$ in these gamma integrals gives
\begin{equation}\label{eq:J-expansion}
 J_q(b)=\frac{\sqrt2\,\Gamma(q+3/2)}{(2b)^{q+3/2}}
 \left[1-\frac{q+3/2}{8b}+O_q(b^{-2})\right].
\end{equation}
Every error assertion follows from the displayed remainder bound and
the exponentially small tail; none extends a finite-$b$ probability
measure by declaration. Since $J_0>0$, division gives
\[
 \mathbb E_b[y]=\frac{J_1}{J_0}=\frac3{4b}+O(b^{-2}),\qquad
 \mathbb E_b[y^2]=\frac{J_2}{J_0}=\frac{15}{16b^2}+O(b^{-3}).
\]
The character $\chi_j$ is the polynomial $U_{d-1}(\cos\theta)$, where
$U_0(x)=1$, $U_1(x)=2x$, and
$U_{r+1}(x)=2xU_r(x)-U_{r-1}(x)$. Induction at $x=1$ gives
\[
 U_{d-1}(1)=d,\qquad U'_{d-1}(1)=\frac{d(d^2-1)}3.
\]
For the derivative formula, differentiate the recurrence; the induction
step is
$U'_{r+1}(1)=2(r+1)+2U'_r(1)-U'_{r-1}(1)$, which the displayed
cubic expression satisfies with its initial values. Polynomial Taylor
expansion on the whole compact interval $y\in[0,2]$ therefore yields
\[
 \frac{\chi_j(\theta)}d
 =1-\frac{d^2-1}{3}y+R_j(y),\qquad |R_j(y)|\leq C_jy^2.
\]
Its exact $\mathbb E_b$ average is $m_j(b)$. Substitute the two moment
estimates and use $(d^2-1)/4=j(j+1)$ to prove \eqref{eq:large-b}.
\end{proof}

Substitution of the retained $b_t$ gives the exact infinitesimal
coefficient
\begin{equation}\label{eq:epsilonmultiplier}
 m_j(b_t(\epsilon))
 =1-\epsilon\frac{2g_{\rm YM}^2}{a}j(j+1)+O_j(\epsilon^2).
\end{equation}
For the chosen basis, $-\Delta^T$ has eigenvalue $j(j+1)$ on spin $j$.
Indeed take the displayed Pauli basis and put $J_a=-i\,dD^j(T_a)$.
Its anti-Hermitian generators satisfy $[T_a,T_b]=-\varepsilon_{abc}T_c$,
so the Hermitian generators satisfy $[J_a,J_b]=i\varepsilon_{abc}J_c$.
Let $J_\pm=J_1\pm iJ_2$. Directly from these commutators,
$[J_3,J_\pm]=\pm J_\pm$ and
$\sum_aJ_a^2=J_-J_++J_3^2+J_3$. On the highest vector
$e_1^{\otimes 2j}$ of the symmetric power, $J_3$ has eigenvalue $j$
and $J_+$ is zero, so the sum of squares has value $j(j+1)$.
Expanding its commutator with each $J_a$ shows it commutes with all
three; irreducibility therefore extends this value to the whole
representation. An orthogonal change to any other specified
orthonormal $T$ basis leaves the sum of squares unchanged.
For $j=1/2$, it is explicitly
$-\sum_aT_a^2=3I/4$, fixing the absolute scale. Thus
\eqref{eq:epsilonmultiplier} has precisely the $\Delta^T$ convention,
and \eqref{eq:metricdictionary} supplies the other metric.

On the full finite edge product define
\begin{equation}\label{eq:freeproduct}
 (K_\epsilon f)(U')=
 \int_{\mathcal Q}\prod_{e\in\mathsf E}
 k_{b_t(\epsilon)}(U'_eU_e^{-1})f(U)d\lambda(U),
 \qquad
 H_0=-\frac{2g_{\rm YM}^2}{a}\sum_e\Delta_e^T.
\end{equation}
A product matrix coefficient with spins $\boldsymbol j=(j_e)$ has
$K_\epsilon$ eigenvalue $\prod_em_{j_e}(b_t(\epsilon))$ and
$H_0$ eigenvalue $(2g_{\rm YM}^2/a)\sum_ej_e(j_e+1)$.
For every finite linear combination $f$ of such coefficients,
\begin{equation}\label{eq:corederivative}
 \lim_{\epsilon\downarrow0}\frac{K_\epsilon f-f}{\epsilon}=-H_0f.
\end{equation}
This follows by the finite product expansion of
\eqref{eq:epsilonmultiplier}, not by an assertion uniform over arbitrary
spins. Contraction and density then imply $K_\epsilon f\to f$ for every
$f\in\mathcal H$: approximate $f$ by a fixed finite combination and
bound the two approximation errors by its distance from $f$.

\section{Exact magnetic factor and temporal-link Gauss integration}

For the ordered spatial plaquette holonomy $U_p$, set
\begin{equation}\label{eq:magnetic}
 V(U)=\frac1{2g_{\rm YM}^2a}
        \sum_{p\in\mathsf P_s}(2-\Rea\tr U_p),\qquad
 M_\epsilon f(U)=e^{-\epsilon V(U)/2}f(U).
\end{equation}
Every magnetic plaquette is present in this sum. It is a bounded
nonnegative real smooth multiplication operator, with
$0\leq V\leq2|\mathsf P_s|/(g_{\rm YM}^2a)$.
In particular $\epsilon V=b_s\sum_p(2-\Rea\tr U_p)$ exactly.

Define the vertex gauge average
\begin{equation}\label{eq:Gauss}
 (P_Gf)(U)=\int_{\SU(2)^{\mathsf V}}f(\alpha_gU)\prod_vdg_v.
\end{equation}
Every pullback in the integrand is unitary. Haar inversion proves
$P_G^*=P_G$, Haar invariance in two successive integrals proves
$P_G^2=P_G$, and another left translation of the integration variable
shows that the range consists exactly of the invariant functions.
Thus $P_G$ is the orthogonal projection onto
$\mathcal H_{\rm inv}$. Equation \eqref{eq:gaugedictionary} proves that
it is also the original-$h$ Gauss average. Gauge invariance of spatial
Wilson traces proves $M_\epsilon P_G=P_GM_\epsilon$.
The identity
\[
 (\alpha_gU')_e(\alpha_gU)_e^{-1}
 =g_{s(e)}U'_eU_e^{-1}g_{s(e)}^{-1}
\]
and centrality of $k_b$ prove that $K_\epsilon$ commutes with $P_G$.

To compute the temporal integral, write $U$ and $U'$ for two adjacent
spatial slices and $q_v$ for the temporal link at vertex $v$ from the
first slice to the second. In the original inverse-transport convention,
the oriented temporal plaquette associated to edge $e:s\to t$ is
\begin{equation}\label{eq:temporalplaquette}
 P_e(U',q,U)=U_eq_t(U'_e)^{-1}q_s^{-1}.
\end{equation}
Its inverse is $q_sU'_eq_t^{-1}U_e^{-1}$, and $\Rea\tr A^{-1}=
\Rea\tr A$ for a unitary matrix. Hence
\begin{equation}\label{eq:temporal-conjugacy}
 \Rea\tr P_e=\Rea\tr[(\alpha_qU')_eU_e^{-1}].
\end{equation}
This equality includes every temporal vertex variable; it makes no
global temporal gauge-fixing assumption. Under independent original
vertex changes $h^{(0)}$ and $h^{(1)}$ on the two slices,
\[
 U_e\mapsto(h_s^{(0)})^{-1}U_eh_t^{(0)},\quad
 U'_e\mapsto(h_s^{(1)})^{-1}U'_eh_t^{(1)},\quad
 q_v\mapsto(h_v^{(0)})^{-1}q_vh_v^{(1)}.
\]
Multiplication in \eqref{eq:temporalplaquette} cancels the intermediate
factors and gives $P_e\mapsto(h_s^{(0)})^{-1}P_eh_s^{(0)}$.
This verifies directly its original gauge law.

Define the raw two-slice Wilson kernel, with symmetric magnetic endpoint
weights, by
\begin{equation}\label{eq:rawkernel}
 \mathcal T_\epsilon^{\rm raw}(U',U)
 =e^{-\epsilon V(U')/2}e^{-\epsilon V(U)/2}
 \int\exp\!\left[-b_t\sum_e(2-\Rea\tr P_e(U',q,U))\right]\prod_vdq_v.
\end{equation}
The integral is over the whole non-Abelian group at every temporal
vertex. With
\begin{equation}\label{eq:scalar}
 c_\epsilon=[e^{-2b_t}Z(b_t)]^N,\qquad
 F_\epsilon=M_\epsilon K_\epsilon M_\epsilon,
\end{equation}
equations \eqref{eq:kernel} and \eqref{eq:temporal-conjugacy} give the
exact operator identity
\begin{equation}\label{eq:fulltransfer}
 \mathcal T_\epsilon^{\rm raw}
 =c_\epsilon M_\epsilon P_GK_\epsilon M_\epsilon
 =c_\epsilon F_\epsilon P_G
 =c_\epsilon P_GF_\epsilon.
\end{equation}
In particular the normalized transfer is $F_\epsilon$ on
$\mathcal H_{\rm inv}$ and $F_\epsilon P_G$ on $\mathcal H$.
The operator $F_\epsilon$ is positive, self-adjoint, injective and a
contraction because $K_\epsilon$ has those properties and
$M_\epsilon$ is positive, bounded, and invertible. Injectivity follows
from
$\langle f,F_\epsilon f\rangle=
\langle M_\epsilon f,K_\epsilon M_\epsilon f\rangle$ and the strictly
positive eigenvalues of $K_\epsilon$. All these statements remain
valid after restriction to the invariant subspace.

For several adjacent slices, the two half magnetic factors at each
interior slice multiply to $e^{-\epsilon V}$. On an open time interval
the two boundary slices carry half weights; on a periodic time lattice
all slices have their full weight. Fubini's theorem therefore identifies
the exact product of \eqref{eq:rawkernel} with the full anisotropic
Wilson integral, retaining all temporal and spatial plaquettes. In the
periodic case with $m$ time steps the exact partition identity is
\begin{equation}\label{eq:partition}
 Z_{\rm Wilson}^{(m)}
 =c_\epsilon^m\tr_{\mathcal H}(F_\epsilon^mP_G)
 =c_\epsilon^m\tr_{\mathcal H_{\rm inv}}(F_\epsilon^m).
\end{equation}
The traces are well defined: the proof of Theorem \ref{thm:multiplier}
gives a trace-class single-link operator, a finite tensor product is
trace class, and multiplication by bounded operators and $P_G$ preserves
that property. The integral and trace equality can also be checked by
successively integrating its positive continuous compact kernels.
Equation \eqref{eq:partition} is an exact finite-$\epsilon$ identity;
the strong convergence below alone is not used to assert trace-norm
convergence of these partition functions.

The Wilson constant and the normalization are explicitly retained as the
common energy shift
\begin{equation}\label{eq:shift}
 E_{\rm scalar}(\epsilon)
 =-\frac1\epsilon\log c_\epsilon
 =\frac N\epsilon\bigl[2b_t-\log Z(b_t)\bigr].
\end{equation}
In particular $0<c_\epsilon<1$, because
$e^{-2b}Z(b)=\int e^{-b(2-\Rea\tr U)}dU$ is a strictly positive average
of a function at most one and strictly below one on a set of positive
Haar measure. From the proved $J_0$ expansion,
\[
 e^{-2b}Z(b)=\frac1{2\sqrt\pi\,b^{3/2}}
             \left[1-\frac3{16b}+O(b^{-2})\right].
\]
Thus the shift is not finite as $\epsilon\downarrow0$; it is a known
scalar that affects all energies equally. The exact relation
\eqref{eq:fulltransfer}, rather than an unrecorded removal of constants,
specifies which operator has the finite time limit.

\section{The actual core and the strong time-continuum limit}

Let $\mathscr D$ be the finite linear span of product matrix coefficients
of all edge representations, and let $\mathscr D_{\rm inv}=P_G\mathscr D$.
Peter--Weyl completeness makes $\mathscr D$ dense in $\mathcal H$ and
$\mathscr D_{\rm inv}$ dense in $\mathcal H_{\rm inv}$. The gauge action
mixes only matrix coefficients with the same finite list of edge spins,
so $P_G\mathscr D\subset\mathscr D$. The diagonal positive operator $H_0$
in \eqref{eq:freeproduct} is self-adjoint on
\[
 D(H_0)=\left\{f:\sum_{\boldsymbol j,\boldsymbol r,\boldsymbol s}
 \left(\frac{2g_{\rm YM}^2}{a}\sum_ej_e(j_e+1)\right)^2
 |f_{\boldsymbol j,\boldsymbol r,\boldsymbol s}|^2<\infty\right\},
\]
where coefficients use an orthonormal matrix-coefficient basis.
Finite spectral truncation proves that $\mathscr D$ is an operator core:
both $f$ and $H_0f$ are approximated by the same finite sums. The
projection $P_G$ commutes with $H_0$, so applying it to these graph-norm
approximations proves the invariant core assertion as well.

Since $V$ in \eqref{eq:magnetic} is bounded and real, the operator
\begin{equation}\label{eq:resultH}
 H=H_0+V
 =-\frac{2g_{\rm YM}^2}{a}\sum_e\Delta_e^T
 +\frac1{2g_{\rm YM}^2a}\sum_{p\in\mathsf P_s}(2-\Rea\tr U_p)
\end{equation}
is self-adjoint and nonnegative on $D(H_0)$. Bounded perturbation does
not change the core, since the graph norms of $H_0$ and $H_0+V$ are
equivalent by the two inequalities
$\|Hf\|\leq\|H_0f\|+\|V\|\|f\|$ and its converse. All gauge
invariance statements persist. The original-metric expression is exactly
\begin{equation}\label{eq:resultHc}
 H=-\frac{g_{\rm YM}^2}{2a}\sum_e\Delta_e^c
 +\frac1{2g_{\rm YM}^2a}\sum_{p\in\mathsf P_s}(2-\Rea\tr U_p).
\end{equation}
No conventional redefinition of $g_{\rm YM}$ was used to obtain either
expression.

For $f\in\mathscr D$, use the exact algebraic decomposition
\[
 (F_\epsilon-I)f=(M_\epsilon-I)f
    +M_\epsilon(K_\epsilon-I)f
    +M_\epsilon K_\epsilon(M_\epsilon-I)f.
\]
Boundedness of $V$ gives
$(M_\epsilon-I)/\epsilon\to-V/2$ in operator norm. Equation
\eqref{eq:corederivative}, the strong convergence $K_\epsilon\to I$,
and the contraction bounds therefore prove
\begin{equation}\label{eq:Fderivative}
 \frac{F_\epsilon f-f}{\epsilon}\longrightarrow-Hf
 \quad\text{for every }f\in\mathscr D,
\end{equation}
and in particular for every $f\in\mathscr D_{\rm inv}$. The magnetic
operator need not commute with the kinetic operator for this argument.

\begin{theorem}\label{thm:stronglimit}
For every fixed $t>0$ and every $f\in\mathcal H$,
\begin{equation}\label{eq:stronglimit}
 F_{t/m}^{\,m}f\longrightarrow e^{-tH}f
 \qquad(m\longrightarrow\infty).
\end{equation}
Consequently the exact normalized Wilson transfer obeys
\begin{equation}\label{eq:projectedlimit}
 \left(c_{t/m}^{-1}\mathcal T_{t/m}^{\rm raw}\right)^mf
 =F_{t/m}^{\,m}P_Gf\longrightarrow e^{-tH}P_Gf.
\end{equation}
On $\mathcal H_{\rm inv}$ the limit is the strongly continuous
semigroup $e^{-tH}$, with generator \eqref{eq:resultH} or equivalently
\eqref{eq:resultHc}. At $t=0$ its identity is $I$ on that subspace;
on the unrestricted space the projected family has value $P_G$.
\end{theorem}
\begin{proof}
We give the positive-contraction specialization of the product-formula
argument directly; the general framework is Chernoff's theorem
\cite{Chernoff}. Put
$A_\epsilon=(I-F_\epsilon)/\epsilon$. It is bounded, self-adjoint and
nonnegative, and \eqref{eq:Fderivative} gives
$A_\epsilon f\to Hf$ on the operator core $\mathscr D$.
For any $\lambda_0>0$ and $f\in\mathscr D$,
\begin{equation}\label{eq:resolventidentity}
 (\lambda_0+A_\epsilon)^{-1}(\lambda_0+H)f-f
 =(\lambda_0+A_\epsilon)^{-1}(H-A_\epsilon)f\longrightarrow0.
\end{equation}
The resolvents have norms at most $1/\lambda_0$. Moreover
$(\lambda_0+H)\mathscr D$ is dense: for arbitrary $u$, apply a
graph-norm approximation from $\mathscr D$ to
$f=(\lambda_0+H)^{-1}u$. Thus the uniform bound extends
\eqref{eq:resolventidentity} to all vectors and proves
$(\lambda_0+A_\epsilon)^{-1}\to(\lambda_0+H)^{-1}$ strongly.

For completeness this implies $e^{-tA_\epsilon}\to e^{-tH}$ strongly
for each fixed $t>0$ without any estimate of an unknown spectrum.
Set $R_\epsilon=(\lambda_0+A_\epsilon)^{-1}$ and
$R=(\lambda_0+H)^{-1}$. On $[0,1/\lambda_0]$ define the continuous
function
\[
 h_t(y)=\begin{cases}\exp[-t(y^{-1}-\lambda_0)],&y>0,\\0,&y=0.
 \end{cases}
\]
Uniform polynomial approximation on this interval and strong convergence
of every polynomial in the uniformly bounded $R_\epsilon$ prove
$h_t(R_\epsilon)\to h_t(R)$ strongly. The spectral calculus identifies
these two operators with the desired semigroups.

It remains to compare powers with this exponential. For integer $m\geq1$
and $0\leq x\leq1$,
\begin{equation}\label{eq:scalarpower}
 0\leq e^{-mx}-(1-x)^m\leq\frac1m.
\end{equation}
For $x\leq1/2$, the integral identity
$-\log(1-x)-x=\int_0^x u/(1-u)du\leq x^2$ gives
\[
 e^{-mx}-(1-x)^m
 \leq mx^2e^{-mx}\leq\frac4{e^2m}<\frac1m.
\]
For $x\geq1/2$ the difference is at most $e^{-m/2}\leq2/(em)<1/m$;
the last inequality follows by maximizing $s e^{-s/2}$ for $s>0$.
The endpoint $x=1$ also satisfies the bound. Spectral calculus with
$x=I-F_\epsilon=\epsilon A_\epsilon$ now yields
\[
 \|F_\epsilon^m-e^{-m\epsilon A_\epsilon}\|\leq1/m.
\]
Take $\epsilon=t/m$ and combine with the established strong convergence
of $e^{-tA_{t/m}}$. This proves \eqref{eq:stronglimit}.
Finally \eqref{eq:fulltransfer}, $P_G^2=P_G$, and commutation with
$F_\epsilon$ prove \eqref{eq:projectedlimit}. The same reasoning
applies directly on the verified invariant core.
\end{proof}

\section{Effective Hamiltonians and the exact finite-vacuum bridge}

On $\mathcal H_{\rm inv}$, positivity and injectivity allow the
self-adjoint functional-calculus definition
\begin{equation}\label{eq:effectiveH}
 H_\epsilon=-\epsilon^{-1}\log(F_\epsilon|_{\mathcal H_{\rm inv}}),
 \qquad
 H_\epsilon^{\rm raw}=E_{\rm scalar}(\epsilon)I+H_\epsilon.
\end{equation}
Zero is not an eigenvalue of $F_\epsilon$, so the logarithm is densely
defined even though zero is a spectral accumulation point. The operator
$H_\epsilon$ is nonnegative. Equation \eqref{eq:effectiveH} is the exact
scalar-energy dictionary, and the raw and normalized effective operators
have identical energy differences and eigenvectors at each $\epsilon$.

They satisfy $H_\epsilon\to H|_{\mathcal H_{\rm inv}}$ in the strong
resolvent sense. Here is a direct consequence of the proof already given,
which also covers arbitrary $\epsilon\downarrow0$ rather than only
$t/m$. Fix $t>0$ and set $m_\epsilon=\lfloor t/\epsilon\rfloor$ and
$t_\epsilon=m_\epsilon\epsilon$. The power comparison gives
$\|F_\epsilon^{m_\epsilon}-e^{-t_\epsilon A_\epsilon}\|
\leq1/m_\epsilon$. For $t_\epsilon\geq t/2$, the scalar mean-value
bound yields
\[
 \sup_{x\geq0}|e^{-t_\epsilon x}-e^{-tx}|
 \leq|t-t_\epsilon|\sup_{x\geq0}xe^{-(t/2)x}
 \leq\frac{2\epsilon}{et}\longrightarrow0.
\]
Thus $F_\epsilon^{m_\epsilon}\to e^{-tH}$ strongly. If
$t/\epsilon=m_\epsilon+\delta$ with $0\leq\delta<1$, then for
$0\leq z\leq1$,
\[
 0\leq z^{m_\epsilon}-z^{m_\epsilon+\delta}
 \leq z^{m_\epsilon}(1-z)
 \leq\frac1{m_\epsilon+1}.
\]
The middle inequality uses $z^\delta\geq z$, and the last follows by
maximizing $z^m(1-z)$. Applying spectral calculus gives
$e^{-tH_\epsilon}=F_\epsilon^{t/\epsilon}\to e^{-tH}$ strongly.
For any $\lambda_0>0$ the resolvent identity
\[
 (\lambda_0+H_\epsilon)^{-1}f
 =\int_0^\infty e^{-\lambda_0t}e^{-tH_\epsilon}f\,dt
\]
and the uniform contraction bound permit dominated convergence in norm,
proving the asserted strong resolvent convergence.

The resulting Hamiltonian is precisely the finite interacting one in
the companion derivation, with the explicit substitution
\begin{equation}\label{eq:companiondictionary}
 n=2,\qquad \kappa_e=\frac{2g_{\rm YM}^2}{a},\qquad
 b_p^{\rm Hamiltonian}=\frac1{2g_{\rm YM}^2a}
                       =\frac{b_s(\epsilon)}\epsilon.
\end{equation}
All magnetic interactions survive this limit. This yields a concrete
state-space and operator map from the finite anisotropic Wilson integral
to the vacuum-weighted finite-lattice identities; it does not identify
the interacting vacuum with Haar measure.

For example, let $\psi_0>0$ be the normalized exact ground state of
\eqref{eq:resultH}, $E_0$ its eigenvalue, and $d\mu=\psi_0^2d\lambda$.
The finite compact elliptic Hamiltonian has this unique ground state
by the positivity and compact-resolvent argument in the companion
paper. For any nonzero smooth gauge-invariant $F$ with $\mu(F)=0$, put
$u=\psi_0F$ and $D=\|u\|^2>0$. The exact ground-state product rule gives
\begin{equation}\label{eq:vacuumform}
 \langle u,(H-E_0)u\rangle
 =\frac{2g_{\rm YM}^2}{a}\sum_{e,a}\int|X_{e,a}F|^2d\mu.
\end{equation}
To see the cancellation, expand
$X^2(\psi_0F)=\psi_0X^2F+2(X\psi_0)XF+F X^2\psi_0$,
subtract $F(H-E_0)\psi_0=0$, and integrate each Haar-skew derivative
by parts. The remaining integral is exactly \eqref{eq:vacuumform}.
Thus no component of $V$ was omitted in passing to that form.

The strong transfer limit proves the concrete correlation identity
\begin{equation}\label{eq:correlation}
 \lim_{m\to\infty}
 \frac{\langle u,F_{t/m}^{\,m}u\rangle}
      {\langle\psi_0,F_{t/m}^{\,m}\psi_0\rangle}
 =\langle u,e^{-t(H-E_0)}u\rangle
 =D\int_{[0,\infty)}e^{-t\omega}d\nu_F(\omega),
\end{equation}
where $\nu_F(A)=D^{-1}\langle u,\mathbf1_A(H-E_0)u\rangle$.
The denominator converges to $e^{-tE_0}>0$. The raw scalar
$c_{t/m}^m$ cancels exactly if both numerator and denominator are
written using the raw transfer. The state $\psi_0$ in
\eqref{eq:correlation} is the actual limiting finite-lattice vacuum,
not a constant trial state. The spectral theorem and
\eqref{eq:vacuumform} give its first moment
\[
 \int\omega\,d\nu_F(\omega)
 =\frac{(2g_{\rm YM}^2/a)\sum_{e,a}\mu(|X_{e,a}F|^2)}{\mu(|F|^2)}.
\]
These are fixed-spatial-regulator statements. Strong time convergence
alone does not supply a spatial continuum vacuum, a limiting spectral
gap, or a gapless interacting four-dimensional quantum theory.

\section{Exact checks and provenance}

The companion script \texttt{check\_temporal\_coefficients.py} completed
237 exact scalar assertions plus exact matrix-entry checks. It verifies
the two anisotropic action coefficients, Bessel recurrence coefficients
through bounded formal orders, the character derivatives for dimensions
1 through 9, the resulting spin coefficients, the SU(2) metric and
Casimir constants, a noncommuting temporal-plaquette example and its
original gauge law, the $g=h^{-1}$ dictionary, and the magnetic endpoint
weights. The all-spin and operator convergence proofs are the analytic
arguments above; they are not inferred from finitely many checked spins.
No numerical integration, spatial continuum calculation, or approximate
vacuum was used.

\begin{thebibliography}{9}
\bibitem{PeterWeyl}
F. Peter and H. Weyl, \emph{Die Vollst\"andigkeit der primitiven
Darstellungen einer geschlossenen kontinuierlichen Gruppe},
Mathematische Annalen \textbf{97} (1927), 737--755,
\href{https://doi.org/10.1007/BF01447892}{doi:10.1007/BF01447892}.

\bibitem{Chernoff}
P. R. Chernoff, \emph{Note on product formulas for operator semigroups},
Journal of Functional Analysis \textbf{2} (1968), 238--242,
\href{https://doi.org/10.1016/0022-1236(68)90020-7}
{doi:10.1016/0022-1236(68)90020-7}.
The positive self-adjoint contraction specialization needed here is
proved directly in Theorem \ref{thm:stronglimit}.
\end{thebibliography}
\end{document}
